% ============================================================================
% §1 Giới thiệu — bám đề cương NCKH (docs/nckh/de_cuong), văn phong tiếng Việt
% học thuật, hạn chế trộn lẫn thuật ngữ tiếng Anh không cần thiết.
% ============================================================================

\section{Giới thiệu}
\label{sec:intro}

Trong bối cảnh chuyển đổi số diễn ra mạnh mẽ trong lĩnh vực giáo dục, việc đo lường và theo dõi năng lực học sinh trở thành một yêu cầu then chốt đối với chất lượng dạy và học. Tuy nhiên, thực tiễn hiện nay cho thấy việc đánh giá năng lực học sinh vẫn chủ yếu dựa trên các kết quả kiểm tra mang tính thời điểm, không phản ánh được bản chất động của quá trình học tập. Trong lý thuyết đo lường giáo dục, cần phân biệt hai quan điểm về năng lực: \emph{năng lực tĩnh} giả định \(\theta_i\) là hằng số cho từng học sinh --- giả định cốt lõi của Lý thuyết Đáp ứng câu hỏi (IRT) truyền thống \cite{lord1980,hambleton1991}; và \emph{năng lực động} \(\theta_{it}\) thay đổi theo thời gian \(t\), phản ánh quá trình học tập liên tục \cite{kim2023,pengwei2022}. Khi học tập diễn ra song song với đánh giá, giả định năng lực tĩnh không còn phù hợp.

IRT đã được chứng minh là công cụ hiệu quả để đo lường năng lực tiềm ẩn thông qua việc mô hình hoá xác suất trả lời đúng \cite{chalmers2012}, với ưu điểm là khả năng tách biệt năng lực người học khỏi các đặc tính của đề thi. Tuy nhiên, hầu hết các mô hình IRT truyền thống đều dựa trên giả định năng lực tĩnh \cite{lord1980}. Hệ thống học trực tuyến OLM (\url{https://olm.vn}) tại Việt Nam hiện sử dụng mô hình Rasch 1PL --- chỉ ước lượng duy nhất một giá trị \(\theta\) cho mỗi học sinh tại một thời điểm và chưa có cơ chế theo dõi sự thay đổi năng lực qua nhiều kỳ thi. Trong khi đó, các mô hình phương trình cấu trúc dọc (Longitudinal Structural Equation Models, LSEM) cho phép theo dõi sự thay đổi của biến tiềm ẩn theo thời gian \cite{mcardle2009}. Khi kết hợp IRT với LSEM, có thể xây dựng được mô hình đánh giá năng lực động phản ánh sát hơn thực tế học tập; tuy vậy, việc khai thác hiệu quả nguồn dữ liệu lớn từ các hệ thống học trực tuyến để phục vụ phân tích quỹ đạo năng lực vẫn là một thách thức chưa được giải quyết đầy đủ.

\textbf{Nghiên cứu ngoài nước.} Hướng nghiên cứu IRT trực tuyến (Online IRT, OIRT) sử dụng Suy diễn Biến phân (Variational Inference) cho phép cập nhật ước lượng năng lực liên tục khi có dữ liệu mới \cite{pengwei2022}. Mô hình VTIRT (Variational Temporal IRT) đã chứng minh khả năng theo dõi năng lực động trên chín bộ dữ liệu thực tế với tốc độ vượt MCMC nhiều bậc độ lớn \cite{kim2023}. Các mô hình DSEM được phát triển để xử lý dữ liệu dọc, cho phép phân tích sự biến đổi năng lực ở cả cấp cá nhân lẫn cấp nhóm \cite{asparouhov2018}; khi khoảng cách giữa các đợt đánh giá không đều, biến thể thời gian liên tục CT-DSEM được sử dụng để bảo toàn tính xác định của mô hình \cite{driver2017}. Mô hình RI-CLPM được dùng để phân tách hiệu ứng đặc trưng ổn định của cá nhân (trait-like) khỏi biến động trong-người tạm thời (within-person) \cite{hamaker2015}. Trên thực tế triển khai, Duolingo sử dụng phương pháp Hồi quy Bán chu kỳ (Half-Life Regression) \cite{settles2016}; Khan Academy dùng Bayesian Knowledge Tracing; IXL dùng đánh giá chẩn đoán dựa trên IRT --- tuy nhiên các hệ thống này chủ yếu theo dõi năng lực trong một phiên học, chưa kết hợp IRT với LSEM để phân tích quỹ đạo tăng trưởng qua nhiều kỳ thi \cite{cambium2020}.

\textbf{Nghiên cứu trong nước.} Tại Việt Nam, ứng dụng IRT chủ yếu dừng lại ở việc đánh giá độ khó đề thi và ước lượng năng lực tĩnh \cite{nguyen_tan}. Hệ thống OLM đã tích luỹ một lượng dữ liệu học tập rất lớn nhưng hiện chỉ sử dụng Rasch 1PL cho ước lượng tĩnh. Việc kết hợp IRT với các mô hình cấu trúc dọc và ước lượng năng lực động trên dữ liệu quy mô lớn vẫn là một khoảng trống nghiên cứu đáng kể, đặc biệt trong bối cảnh dữ liệu giáo dục Việt Nam.

\textbf{Khoảng trống và phát biểu bài toán.} Từ phân tích trên, ba khoảng trống cụ thể được nhận diện: (i) thiếu các mô hình tích hợp IRT--LSEM phục vụ đo lường năng lực động trong bối cảnh dữ liệu giáo dục Việt Nam; (ii) chưa có thuật toán ước lượng năng lực có khả năng cập nhật đệ quy sau mỗi đợt đánh giá phù hợp với dữ liệu lớn từ hệ thống học trực tuyến; (iii) chưa khai thác dữ liệu thực tế để phân tích quỹ đạo phát triển năng lực và hỗ trợ ra quyết định sư phạm. Bài toán nghiên cứu được phát biểu như sau: \emph{cho ma trận đáp ứng \(Y_{ijt} \in \{0,1\}\) của \(N\) học sinh, \(J_t\) câu hỏi qua \(T\) đợt đánh giá, hãy đồng thời ước lượng tham số câu hỏi \(b_j\), năng lực động \(\theta_{i1}, \theta_{i2}, \dots, \theta_{iT}\) tại mỗi thời điểm và mô hình hoá quỹ đạo thay đổi của chuỗi năng lực này, đồng thời cung cấp cơ chế cập nhật đệ quy khi có dữ liệu mới.}

\textbf{Đóng góp của bài báo.} (1) Đề xuất một khung tích hợp IRT--LSEM tổng quát kèm pipeline có khả năng tái lập đầy đủ (B0--B7), bao gồm tầng đo lường IRT 1PL với chấm điểm Bayes EAP và bốn biến thể cấu trúc (LGCM, LCSM, mô hình tăng trưởng đa cấp, CT-DSEM thời gian liên tục cho khoảng cách thời gian không đều). (2) Cung cấp bằng chứng định lượng từ \(N = 24{.}213\) học sinh THPT Việt Nam cho thấy mô hình động vượt mô hình tĩnh cơ sở với \(\Delta\)AIC = 82--167 và \(p\)-value LRT \(\le 10^{-19}\) --- đủ căn cứ thống kê để bác bỏ giả định năng lực tĩnh ở quy mô triển khai thực tế. (3) Bộ lọc Kalman cục bộ sử dụng phương sai quan sát biến thiên theo thời gian (lấy từ sai số chuẩn EAP) giúp giảm 32--41\% phương sai sai số đo lường \emph{mà không cần thêm câu hỏi khảo sát}. (4) Hai ứng dụng kế tiếp --- cảnh báo sớm học sinh có dấu hiệu giảm năng lực và gợi ý câu hỏi tối đa hoá thông tin --- được trình diễn trên quỹ đạo thực tế.

\textbf{Cấu trúc bài báo.} Mục~\ref{sec:content} trình bày phương pháp, dữ liệu, kết quả và thảo luận; Mục~\ref{sec:conclusion} kết luận và đề xuất hướng phát triển.
